\section{Introduction}\label{sec:introduction}

For complex functions $F,G\in L^3(\R^2)$ and
$0<\varepsilon<R<\infty$, define the annular triangular Hilbert transform by
\begin{equation}\label{eq:bilinear-truncation}
 B_{\varepsilon,R}(F,G)(x,y)
 =\int_{\varepsilon<|t|<R}F(x+t,y)G(x,y+t)\,\frac{\dd t}{t}.
\end{equation}
The two inputs are translated along different coordinate directions.
Every finite integral is absolutely convergent outside one common
null set; we use its usual representative there and set all quantities
to zero on the exceptional set. Lemma~\ref{lem:finite-truncations}
justifies this convention and continuity in the endpoints.

For $r>2$, its annular variation is
\begin{equation}\label{eq:variation-definition}
 V_r(F,G)(x,y)
 =\sup_{\substack{J\ge1\,,\ 0<t_0<\cdots<t_J\\t_j\in\mathbb Q}}
    \left(\sum_{j=1}^J
          |B_{t_{j-1},t_j}(F,G)(x,y)|^r\right)^{1/r}.
\end{equation}
The supremum is pointwise: the partition may depend on $(x,y)$, and
there is no restriction on the number of endpoints within a dyadic
scale interval. The use of rational endpoints makes measurability
immediate and, by endpoint continuity, does not change the supremum.

\begin{theorem}[Full annular variation]\label{thm:variation}
For every real $r>2$ there is a finite constant $C_r$ such that
\[
 \|V_r(F,G)\|_{L^{3/2}(\R^2)}
 \le C_r\|F\|_{L^3(\R^2)}\|G\|_{L^3(\R^2)}
\]
for all complex $F,G\in L^3(\R^2)$.
\end{theorem}

Taking a single increment gives the maximal estimate with the supremum
over both hard truncation endpoints. In Section~\ref{sec:consequences}
we prove that the joint principal value
\[
 B(F,G)=\lim_{\substack{\varepsilon\downarrow0\\R\uparrow\infty}}
                   B_{\varepsilon,R}(F,G)
\]
exists almost everywhere and in $L^{3/2}$. In fact, the supremum of
$|B_{\varepsilon,R}(F,G)-B(F,G)|$ over
$0<\varepsilon<1/m$, $R>m$, tends to zero in $L^{3/2}$.

Pairing the bilinear output with a third $L^3$ input gives a scalar
form. Its symmetric, or simplex, coordinates are
\begin{equation}\label{eq:intro-simplex}
 \Lambda_{\varepsilon,R}(G_0,G_1,G_2)
 =\iiint_{\varepsilon<|a+b+c|<R}
  G_0(a,b)G_1(b,c)G_2(c,a)
             \frac{\dd a\dd b\dd c}{a+b+c}.
\end{equation}
The three equal exponents $3$ form the symmetric point of the scalar
H\"older relation $1/p_0+1/p_1+1/p_2=1$.
We obtain a uniform bound by $C\prod_v\|G_v\|_3$ and the joint scalar
principal value for every complex $L^3$ triple.
Proposition~\ref{prop:flat-equivalence} gives the norm-preserving
coordinate maps and equality of the optimal constants for the flat
and simplex scalar formulations. In particular, this proves the
symmetric scalar estimate in Thiele's Problem~13 in
\cite[Section~9]{GrafakosEtAl2017}.

\subsection{History and related work}

The flat operator occurs in Demeter and Thiele's study of bilinear
averages for two commuting transformations
\cite[Section~6]{DemeterThiele2010}. The scalar form
\eqref{eq:intro-simplex} is a basic entangled singular integral: each
input sees one pair of the three variables, and the pairs form a cycle.
This dependence differs from the bipartite configurations in many
multilinear estimates. Kova\v{c}, Thiele and Zorin-Kranich established
estimates for a Walsh model with two general inputs and a third input
of specified structure \cite[Theorem~1.7]{KovacThieleZorinKranich2015}.
Their paper also records the cyclic trace formulation and the
equivalence of the flat and simplex forms.

Quantitative cancellation over a finite range of scales was developed
in several stages. Tao proved sublogarithmic bounds for one-dimensional
multilinear Hilbert transforms \cite{Tao2016Cancellation}, and
Zorin-Kranich obtained sublogarithmic cancellation for simplex Hilbert
transforms \cite{ZorinKranich2017}. Durcik, Kova\v{c}, \v{S}kreb and
Thiele proved square-function and norm-variation estimates that yield
square-root cancellation for the triangular form over finitely many
smooth scales \cite[Corollary~5]{DurcikKovacSkrebThiele2019}.
Durcik, Kova\v{c} and Thiele obtained power-type cancellation for hard
truncations \cite[Theorem~1]{DurcikKovacThiele2019}. Permuting and
interpolating their $(4,4,2)$ bounds gives a
$\sqrt{\log(R/\varepsilon)}$ bound at the symmetric tuple.
These estimates retain a dependence on the scale range. Moreover,
norm-variation estimates for a fixed list of scales do not place the
supremum over output-dependent partitions inside the output norm as
in \eqref{eq:variation-definition}.

Other results gain cancellation by averaging a direction parameter
\cite{DurcikRoos2021,LinSlavikova2026} or by introducing curvature
\cite{ChristDurcikRoos2021,HsuLin2026}. Those mechanisms are distinct
from the flat problem considered here. The broader role of variation
in strengthening maximal estimates and controlling convergence is
illustrated by L\'epingle's martingale inequality
\cite{Lepingle1976}, the variational inequalities of Jones, Seeger and
Wright \cite{JonesSeegerWright2008}, and the variation-norm Carleson
theorem \cite{OberlinSeegerTaoThieleWright2012}. These results provide
context rather than an input to our proof.

The energy method here is related to telescoping arguments for
entangled forms: Kova\v{c}'s twisted paraproduct and Bellman-function
framework \cite{Kovac2012Twisted,Kovac2011}, and Durcik's continuous
Gaussian telescoping \cite{Durcik2015,Durcik2017}. The obstruction
presented by a triangular cycle is discussed in
\cite[Section~6.2]{Kovac2011}. We work with the finite matrix energy
$\tr(P^{3/2})$ for $P\ge0$, establish a mixed trace inequality, and compare
its heat derivatives at neighboring vertices. All matrix and heat
estimates required below are proved in this paper.

\subsection{The count estimate and the new ingredients}

The proof reduces variation to a quantitative estimate for finitely
many disjoint annuli. For an integer $n\ge1$, define
\begin{equation}\label{eq:count-definition}
 \mathcal S_n(F,G)
 =\sup\left\{\sum_{j=1}^m|B_{\varepsilon_j,R_j}(F,G)|:
 \begin{array}{l}
 0\le m\le n,\quad 0<\varepsilon_j<R_j,\quad
     \varepsilon_j,R_j\in\mathbb Q,\\
 (\varepsilon_j,R_j)\text{ pairwise disjoint}
 \end{array}\right\}.
\end{equation}
Intervals may share endpoints, and the empty sum is zero. As in the
variation, the entire choice is made separately at every output point.

\begin{theorem}[Disjoint-annulus count estimate]\label{thm:count}
There is an absolute constant $C$ such that, for every integer $n\ge1$
and every complex $F,G\in L^3(\R^2)$,
\begin{equation}\label{eq:count-bound}
 \|\mathcal S_n(F,G)\|_{3/2}
 \le C\sqrt n\log(2+n)\|F\|_3\|G\|_3.
\end{equation}
\end{theorem}

Theorem~\ref{thm:variation} follows by arranging the increments of
each partition in decreasing order and grouping their ranks dyadically.
The resulting series has terms bounded by
$C(1+j)2^{j(1/r-1/2)}\|F\|_3\|G\|_3$.
The substance of the paper is therefore \eqref{eq:count-bound}.

For smooth annuli, we write the kernel as a scale integral of
convolutions of three one-variable Gaussian derivative windows.
Each convolution is an integral over window centers whose sum is zero;
with the centers fixed, the three window factors depend separately on
$a,b,c$. On a finite grid the resulting sums are cyclic matrix traces:
the sampled pairwise inputs supply three matrices, and the window
factors supply diagonal insertions. We localize the matrices with
positive Gaussian weights. At each vertex we place the two incident
weighted matrices side by side to form a row matrix $R$.
The energy $\tr((R^*R)^{3/2})$ has the same degree-three scaling as the
trilinear form. A mixed trace inequality bounds the localized trace by
quadratic costs from the Hessians of these energies. With the matrices
fixed, heat identities control the accumulated
costs by initial energies after integration over centers. Those
integrated energies are bounded by sums of cubed $L^3$ norms.

There are two different costs in passing from one annulus to many.
For smooth annuli, entries of the matrix formed from the dual input
switch on and off at their chosen endpoints. We estimate the energy jumps entrywise,
then rescale the dual input to obtain a $\sqrt n$ bound.
For hard annuli, estimating each endpoint error separately would
give a linear loss. We instead group endpoints in dyadic radius
intervals. Disjointness gives uniform size and derivative bounds
for each group even when it has many jumps.

The main additional analytic tool is the frequency-block estimate
of Proposition~\ref{prop:frequency-block}. It handles modulated
Gaussian windows with coefficients depending jointly on scale,
frequency, and output position. For uniformly bounded coefficients,
a pointwise square-integral bound over scales and frequencies suffices,
uniformly in the frequency size.
To obtain that uniformity, we narrow one Gaussian variance and
average its center shifts. This expresses each oscillatory insertion
as an average of derivatives of narrower Gaussians. Auxiliary heat
flows retaining only one input edge control the extra derivative
terms. A two-dimensional Gaussian estimate then bounds the energy
changes at the ends of the scale intervals without a loss from the
narrow variance.
This estimate and the rough-kernel family bound in
Proposition~\ref{prop:rough-family} are formulated separately from
the annular application.

\subsection{Organization}

Section~\ref{sec:preliminaries} fixes truncations, coordinates, and
line maximal estimates. Section~\ref{sec:matrix} proves the finite
matrix inequalities; Section~\ref{sec:heat} develops their Gaussian
heat flow. Section~\ref{sec:smooth} proves the repeated-mask estimate
for smooth annuli. Sections~\ref{sec:frequency} and~\ref{sec:rough}
establish the uniform frequency-block estimate and control the hard
endpoint errors. Section~\ref{sec:completion} proves
Theorems~\ref{thm:count} and~\ref{thm:variation} by linearization,
density, and rank summation. Section~\ref{sec:consequences} derives
the maximal, principal-value, and scalar conclusions.
